\documentclass[aspectratio=169]{beamer}

%\documentclass[handout]{beamer}
%% To make 4 per page
%\usepackage{pgfpages}
%\mode<handout>{\setbeamercolor{background canvas}{bg=white}}
%\pgfpagesuselayout{4 on 1}[letterpaper,landscape]%,border shrink=5mm]

\usetheme{default}
% Slide setup, colour independent

\usepackage{amsmath,amssymb,amsthm}
\usepackage{colortbl}
\usepackage{bm}
\usepackage{xcolor}
\usepackage{dsfont}
\usepackage{setspace}
%\usepackage{subfigure}
% To use \ding{234} and the like
\usepackage{pifont}
% To cross reference between slide files
\usepackage{zref-xr,zref-user}
% Use something like
% \zexternaldocument{fileI}
% in the tex files. And cite using \zref instead of \ref

% Fields and the like
\def\IC{\mathbb{C}}
\def\IF{\mathbb{F}}
\def\II{\mathbb{I}}
\def\IJ{\mathbb{J}}
\def\IM{\mathbb{M}}
\def\IN{\mathbb{N}}
\def\IP{\mathbb{P}}
\def\IR{\mathbb{R}}
\def\IZ{\mathbb{Z}}
\def\11{\mathds{1}}


% Bold lowercase
\def\ba{\mathbf{a}}
\def\bb{\mathbf{b}}
\def\bc{\mathbf{c}}
\def\bd{\mathbf{d}}
\def\be{\mathbf{e}}
\def\bf{\mathbf{f}}
\def\bh{\mathbf{h}}
\def\bi{\mathbf{i}}
\def\bj{\mathbf{j}}
\def\bk{\mathbf{k}}
\def\bn{\mathbf{n}}
\def\bp{\mathbf{p}}
\def\br{\mathbf{r}}
\def\bs{\mathbf{s}}
\def\bu{\mathbf{u}}
\def\bv{\mathbf{v}}
\def\bw{\mathbf{w}}
\def\bx{\mathbf{x}}
\def\by{\mathbf{y}}
\def\bz{\mathbf{z}}

% Bold capitals
\def\bB{\mathbf{B}}
\def\bD{\mathbf{D}}
\def\bF{\mathbf{F}}
\def\bG{\mathbf{G}}
\def\bI{\mathbf{I}}
\def\bL{\mathbf{L}}
\def\bN{\mathbf{N}}
\def\bP{\mathbf{P}}
\def\bR{\mathbf{R}}
\def\bS{\mathbf{S}}
\def\bT{\mathbf{T}}
\def\bX{\mathbf{X}}

% Bold numbers
\def\b0{\mathbf{0}}

% Bold greek
\bmdefine{\bmu}{\bm{\mu}}
\def\bphi{\bm{\phi}}
\def\bvarphi{\bm{\varphi}}

% Bold red sentence
\def\boldred#1{{\color{red}\textbf{#1}}}
\def\defword#1{{\color{orange}\textbf{#1}}}

% Caligraphic letters
\def\A{\mathcal{A}}
\def\B{\mathcal{B}}
\def\C{\mathcal{C}}
\def\D{\mathcal{D}}
\def\E{\mathcal{E}}
\def\F{\mathcal{F}}
\def\G{\mathcal{G}}
\def\H{\mathcal{H}}
\def\I{\mathcal{I}}
\def\L{\mathcal{L}}
\def\M{\mathcal{M}}
\def\N{\mathcal{N}}
\def\P{\mathcal{P}}
\def\R{\mathcal{R}}
\def\S{\mathcal{S}}
\def\T{\mathcal{T}}
\def\U{\mathcal{U}}
\def\V{\mathcal{V}}

% tt font for code
\def\code#1{{\tt #1}}

% i.e., e.g.
\def\eg{\emph{e.g.}}
\def\ie{\emph{i.e.}}


% Operators and special symbols
\def\nbOne{{\mathchoice {\rm 1\mskip-4mu l} {\rm 1\mskip-4mu l}
{\rm 1\mskip-4.5mu l} {\rm 1\mskip-5mu l}}}
\def\cov{\ensuremath{\mathsf{cov}}}
\def\Var{\ensuremath{\mathsf{Var}\ }}
\def\Im{\textrm{Im}\;}
\def\Re{\textrm{Re}\;}
\def\det{\ensuremath{\mathsf{det}}}
\def\diag{\ensuremath{\mathsf{diag}}}
\def\nullspace{\ensuremath{\mathsf{null}}}
\def\nullity{\ensuremath{\mathsf{nullity}}}
\def\rank{\ensuremath{\mathsf{rank}}}
\def\range{\ensuremath{\mathsf{range}}}
\def\sgn{\ensuremath{\mathsf{sgn}}}
\def\Span{\ensuremath{\mathsf{span}}}
\def\tr{\ensuremath{\mathsf{tr}}}
\def\imply{$\Rightarrow$}
\def\restrictTo#1#2{\left.#1\right|_{#2}}
\newcommand{\parallelsum}{\mathbin{\!/\mkern-5mu/\!}}
\def\dsum{\mathop{\displaystyle \sum }}%
\def\dind#1#2{_{\substack{#1\\ #2}}}

\DeclareMathOperator{\GL}{GL}
\DeclareMathOperator{\Rel}{Re}
\def\Nt#1{\left|\!\left|\!\left|#1\right|\!\right|\!\right|}
\newcommand{\tripbar}{|\! |\! |}



% The beamer bullet (in base colour)
\def\bbullet{\leavevmode\usebeamertemplate{itemize item}\ }

% Theorems and the like
\newtheorem{proposition}[theorem]{Proposition}
\newtheorem{property}[theorem]{Property}
\newtheorem{importantproperty}[theorem]{Property}
\newtheorem{importanttheorem}[theorem]{Theorem}
%\newtheorem{lemma}[theorem]{Lemma}
%\newtheorem{corollary}[theorem]{Corollary}
\newtheorem{remark}[theorem]{Remark}
\setbeamertemplate{theorems}[numbered]
%\setbeamertemplate{theorems}[ams style]

%
%\usecolortheme{orchid}
%\usecolortheme{orchid}

\def\red{\color[rgb]{1,0,0}}
\def\blue{\color[rgb]{0,0,1}}
\def\green{\color[rgb]{0,1,0}}


% Get rid of navigation stuff
\setbeamertemplate{navigation symbols}{}

% Set footline/header line
\setbeamertemplate{footline}
{%
\quad p. \insertpagenumber \quad--\quad \insertsection\vskip2pt
}
% \setbeamertemplate{headline}
% {%
% \quad\insertsection\hfill p. \insertpagenumber\quad\mbox{}\vskip2pt
% }


\makeatletter
\newlength\beamerleftmargin
\setlength\beamerleftmargin{\Gm@lmargin}
\makeatother

% Colours for special pages
\def\extraContent{yellow!20}


%%%%%%%%%%%%%%%%%
\usepackage{tikz}
\usetikzlibrary{shapes,arrows}
\usetikzlibrary{positioning}
\usetikzlibrary{shapes.symbols,shapes.callouts,patterns}
\usetikzlibrary{calc,fit}
\usetikzlibrary{backgrounds}
\usetikzlibrary{decorations.pathmorphing,fit,petri}
\usetikzlibrary{automata}
\usetikzlibrary{fadings}
\usetikzlibrary{patterns,hobby}

\usetikzlibrary{backgrounds,fit,petri}


\usepackage{pgfplots}
\pgfplotsset{compat=1.6}
\pgfplotsset{ticks=none}

\usetikzlibrary{decorations.markings}
\usetikzlibrary{arrows.meta}
\tikzset{>=stealth}

% For tikz
\usetikzlibrary{shapes,arrows}
\usetikzlibrary{positioning}
\tikzstyle{cloud} = [draw, ellipse,fill=red!20, node distance=0.87cm,
minimum height=2em]
\tikzstyle{line} = [draw, -latex']


%%% For max frame images
\newenvironment{changemargin}[2]{%
\begin{list}{}{%
\setlength{\topsep}{0pt}%
\setlength{\leftmargin}{#1}%
\setlength{\rightmargin}{#2}%
\setlength{\listparindent}{\parindent}%
\setlength{\itemindent}{\parindent}%
\setlength{\parsep}{\parskip}%
}%
\item[]}{\end{list}}


% Make one image take up the entire slide content area in beamer,.:
% centered/centred full-screen image, with title:
% This uses the whole screen except for the 1cm border around it
% all. 128x96mm
\newcommand{\titledFrameImage}[2]{
\begin{frame}{#1}
%\begin{changemargin}{-1cm}{-1cm}
\begin{center}
\includegraphics[width=108mm,height=\textheight,keepaspectratio]{#2}
\end{center}
%\end{changemargin}
\end{frame}
}

% Make one image take up the entire slide content area in beamer.:
% centered/centred full-screen image, no title:
% This uses the whole screen except for the 1cm border around it
% all. 128x96mm
\newcommand{\plainFrameImage}[1]{
\begin{frame}[plain]
%\begin{changemargin}{-1cm}{-1cm}
\begin{center}
\includegraphics[width=108mm,height=76mm,keepaspectratio]{#1}
\end{center}
%\end{changemargin}
\end{frame}
}

% Make one image take up the entire slide area, including borders, in beamer.:
% centered/centred full-screen image, no title:
% This uses the entire whole screen
\newcommand{\maxFrameImage}[1]{
\begin{frame}[plain]
\begin{changemargin}{-1cm}{-1cm}
\begin{center}
\includegraphics[width=\paperwidth,height=\paperheight,keepaspectratio]
{#1}
\end{center}
\end{changemargin}
\end{frame}
}

% This uses the entire whole screen (to include in frame)
\newcommand{\maxFrameImageNoFrame}[1]{
\begin{changemargin}{-1cm}{-1cm}
\begin{center}
\includegraphics[width=\paperwidth,height=0.99\paperheight,keepaspectratio]
{#1}
\end{center}
\end{changemargin}
}

% Make one image take up the entire slide area, including borders, in beamer.:
% centered/centred full-screen image, no title:
% This uses the entire whole screen
\newcommand{\maxFrameImageColor}[2]{
\begin{frame}[plain]
\setbeamercolor{normal text}{bg=#2!20}
\begin{changemargin}{-1cm}{-1cm}
\begin{center}
\includegraphics[width=\paperwidth,height=\paperheight,keepaspectratio]
{#1}
\end{center}
\end{changemargin}
\end{frame}
}


\usepackage{tikz}
\usetikzlibrary{patterns,hobby}
\usepackage{pgfplots}
\pgfplotsset{compat=1.6}
\pgfplotsset{ticks=none}

\usetikzlibrary{backgrounds}
\usetikzlibrary{decorations.markings}
\usetikzlibrary{arrows.meta}
\tikzset{>=stealth}

\tikzset{
  clockwise arrows/.style={
    postaction={
      decorate,
      decoration={
        markings,
        mark=between positions 0.1 and 0.9 step 40pt with {\arrow{>}},
   }}}}


   %%%%%%%%%%%
% To have links to parts in the outline
\makeatletter
\AtBeginPart{%
  \addtocontents{toc}{\protect\beamer@partintoc{\the\c@part}{\beamer@partnameshort}{\the\c@page}}%
}
%% number, shortname, page.
\providecommand\beamer@partintoc[3]{%
  \ifnum\c@tocdepth=-1\relax
    % requesting onlyparts.
    \makebox[6em]{Part #1:} \textcolor{green!30!blue}{\hyperlink{#2}{#2}}
    \par
  \fi
}
\define@key{beamertoc}{onlyparts}[]{%
  \c@tocdepth=-1\relax
}
\makeatother%

\newcommand{\nameofthepart}{}
\newcommand{\nupart}[1]%
    {   \part{#1}%
        \renewcommand{\nameofthepart}{#1}%
        {
          \setbeamercolor{background canvas}{bg=orange!50}
          \begin{frame}{#1}%\partpage 
          \hypertarget{\nameofthepart}{}\tableofcontents%
          \end{frame}
        }
    }



\usecolortheme{orchid}

%% Listings
\usepackage{listings}
\definecolor{mygreen}{rgb}{0,0.6,0}
\definecolor{mygray}{rgb}{0.5,0.5,0.5}
\definecolor{mymauve}{rgb}{0.58,0,0.82}
\definecolor{mygold}{rgb}{1,0.843,0}
\definecolor{myblue}{rgb}{0.537,0.812,0.941}

\definecolor{lgreen}{rgb}{0.6,0.9,.6}
\definecolor{lred}{rgb}{1,0.5,.5}

\lstloadlanguages{R}
\lstset{ %
  language=R,
  backgroundcolor=\color{black!05},   % choose the background color
  basicstyle=\footnotesize\ttfamily,        % size of fonts used for the code
  breaklines=true,                 % automatic line breaking only at whitespace
  captionpos=b,                    % sets the caption-position to bottom
  commentstyle=\color{mygreen},    % comment style
  escapeinside={\%*}{*)},          % if you want to add LaTeX within your code
  keywordstyle=\color{red},       % keyword style
  stringstyle=\color{mygold},     % string literal style
  keepspaces=true,
  columns=fullflexible,
  tabsize=4,
}
% Could also do (in lstset)
% basicstyle==\fontfamily{pcr}\footnotesize
\lstdefinelanguage{Renhanced}%
  {keywords={abbreviate,abline,abs,acos,acosh,action,add1,add,%
      aggregate,alias,Alias,alist,all,anova,any,aov,aperm,append,apply,%
      approx,approxfun,apropos,Arg,args,array,arrows,as,asin,asinh,%
      atan,atan2,atanh,attach,attr,attributes,autoload,autoloader,ave,%
      axis,backsolve,barplot,basename,besselI,besselJ,besselK,besselY,%
      beta,binomial,body,box,boxplot,break,browser,bug,builtins,bxp,by,%
      c,C,call,Call,case,cat,category,cbind,ceiling,character,char,%
      charmatch,check,chol,chol2inv,choose,chull,class,close,cm,codes,%
      coef,coefficients,co,col,colnames,colors,colours,commandArgs,%
      comment,complete,complex,conflicts,Conj,contents,contour,%
      contrasts,contr,control,helmert,contrib,convolve,cooks,coords,%
      distance,coplot,cor,cos,cosh,count,fields,cov,covratio,wt,CRAN,%
      create,crossprod,cummax,cummin,cumprod,cumsum,curve,cut,cycle,D,%
      data,dataentry,date,dbeta,dbinom,dcauchy,dchisq,de,debug,%
      debugger,Defunct,default,delay,delete,deltat,demo,de,density,%
      deparse,dependencies,Deprecated,deriv,description,detach,%
      dev2bitmap,dev,cur,deviance,off,prev,,dexp,df,dfbetas,dffits,%
      dgamma,dgeom,dget,dhyper,diag,diff,digamma,dim,dimnames,dir,%
      dirname,dlnorm,dlogis,dnbinom,dnchisq,dnorm,do,dotplot,double,%
      download,dpois,dput,drop,drop1,dsignrank,dt,dummy,dump,dunif,%
      duplicated,dweibull,dwilcox,dyn,edit,eff,effects,eigen,else,%
      emacs,end,environment,env,erase,eval,equal,evalq,example,exists,%
      exit,exp,expand,expression,External,extract,extractAIC,factor,%
      fail,family,fft,file,filled,find,fitted,fivenum,fix,floor,for,%
      For,formals,format,formatC,formula,Fortran,forwardsolve,frame,%
      frequency,ftable,ftable2table,function,gamma,Gamma,gammaCody,%
      gaussian,gc,gcinfo,gctorture,get,getenv,geterrmessage,getOption,%
      getwd,gl,glm,globalenv,gnome,GNOME,graphics,gray,grep,grey,grid,%
      gsub,hasTsp,hat,heat,help,hist,home,hsv,httpclient,I,identify,if,%
      ifelse,Im,image,\%in\%,index,influence,measures,inherits,install,%
      installed,integer,interaction,interactive,Internal,intersect,%
      inverse,invisible,IQR,is,jitter,kappa,kronecker,labels,lapply,%
      layout,lbeta,lchoose,lcm,legend,length,levels,lgamma,library,%
      licence,license,lines,list,lm,load,local,locator,log,log10,log1p,%
      log2,logical,loglin,lower,lowess,ls,lsfit,lsf,ls,machine,Machine,%
      mad,mahalanobis,make,link,margin,match,Math,matlines,mat,matplot,%
      matpoints,matrix,max,mean,median,memory,menu,merge,methods,min,%
      missing,Mod,mode,model,response,mosaicplot,mtext,mvfft,na,nan,%
      names,omit,nargs,nchar,ncol,NCOL,new,next,NextMethod,nextn,%
      nlevels,nlm,noquote,NotYetImplemented,NotYetUsed,nrow,NROW,null,%
      numeric,\%o\%,objects,offset,old,on,Ops,optim,optimise,optimize,%
      options,or,order,ordered,outer,package,packages,page,pairlist,%
      pairs,palette,panel,par,parent,parse,paste,path,pbeta,pbinom,%
      pcauchy,pchisq,pentagamma,persp,pexp,pf,pgamma,pgeom,phyper,pico,%
      pictex,piechart,Platform,plnorm,plogis,plot,pmatch,pmax,pmin,%
      pnbinom,pnchisq,pnorm,points,poisson,poly,polygon,polyroot,pos,%
      postscript,power,ppoints,ppois,predict,preplot,pretty,Primitive,%
      print,prmatrix,proc,prod,profile,proj,prompt,prop,provide,%
      psignrank,ps,pt,ptukey,punif,pweibull,pwilcox,q,qbeta,qbinom,%
      qcauchy,qchisq,qexp,qf,qgamma,qgeom,qhyper,qlnorm,qlogis,qnbinom,%
      qnchisq,qnorm,qpois,qqline,qqnorm,qqplot,qr,Q,qty,qy,qsignrank,%
      qt,qtukey,quantile,quasi,quit,qunif,quote,qweibull,qwilcox,%
      rainbow,range,rank,rbeta,rbind,rbinom,rcauchy,rchisq,Re,read,csv,%
      csv2,fwf,readline,socket,real,Recall,rect,reformulate,regexpr,%
      relevel,remove,rep,repeat,replace,replications,report,require,%
      resid,residuals,restart,return,rev,rexp,rf,rgamma,rgb,rgeom,R,%
      rhyper,rle,rlnorm,rlogis,rm,rnbinom,RNGkind,rnorm,round,row,%
      rownames,rowsum,rpois,rsignrank,rstandard,rstudent,rt,rug,runif,%
      rweibull,rwilcox,sample,sapply,save,scale,scan,scan,screen,sd,se,%
      search,searchpaths,segments,seq,sequence,setdiff,setequal,set,%
      setwd,show,sign,signif,sin,single,sinh,sink,solve,sort,source,%
      spline,splinefun,split,sqrt,stars,start,stat,stem,step,stop,%
      storage,strstrheight,stripplot,strsplit,structure,strwidth,sub,%
      subset,substitute,substr,substring,sum,summary,sunflowerplot,svd,%
      sweep,switch,symbol,symbols,symnum,sys,status,system,t,table,%
      tabulate,tan,tanh,tapply,tempfile,terms,terrain,tetragamma,text,%
      time,title,topo,trace,traceback,transform,tri,trigamma,trunc,try,%
      ts,tsp,typeof,unclass,undebug,undoc,union,unique,uniroot,unix,%
      unlink,unlist,unname,untrace,update,upper,url,UseMethod,var,%
      variable,vector,Version,vi,warning,warnings,weighted,weights,%
      which,while,window,write,\%x\%,x11,X11,xedit,xemacs,xinch,xor,%
      xpdrows,xy,xyinch,yinch,zapsmall,zip},%
   otherkeywords={!,!=,~,$,*,\%,\&,\%/\%,\%*\%,\%\%,<-,<<-,_,/},%
   alsoother={._$},%
   sensitive,%
   morecomment=[l]\#,%
   morestring=[d]",%
   morestring=[d]'% 2001 Robert Denham
  }%

%%%%%%% 
%% Definitions in yellow boxes
\usepackage{etoolbox}
\setbeamercolor{block title}{use=structure,fg=structure.fg,bg=structure.fg!40!bg}
\setbeamercolor{block body}{parent=normal text,use=block title,bg=block title.bg!20!bg}

\BeforeBeginEnvironment{definition}{%
	\setbeamercolor{block title}{fg=black,bg=yellow!20!white}
	\setbeamercolor{block body}{fg=black, bg=yellow!05!white}
}
\AfterEndEnvironment{definition}{
	\setbeamercolor{block title}{use=structure,fg=structure.fg,bg=structure.fg!20!bg}
	\setbeamercolor{block body}{parent=normal text,use=block title,bg=block title.bg!50!bg, fg=black}
}
\BeforeBeginEnvironment{importanttheorem}{%
	\setbeamercolor{block title}{fg=black,bg=red!20!white}
	\setbeamercolor{block body}{fg=black, bg=red!05!white}
}
\AfterEndEnvironment{importanttheorem}{
	\setbeamercolor{block title}{use=structure,fg=structure.fg,bg=structure.fg!20!bg}
	\setbeamercolor{block body}{parent=normal text,use=block title,bg=block title.bg!50!bg, fg=black}
}
\BeforeBeginEnvironment{importantproperty}{%
	\setbeamercolor{block title}{fg=black,bg=red!50!white}
	\setbeamercolor{block body}{fg=black, bg=red!30!white}
}
\AfterEndEnvironment{importantproperty}{
	\setbeamercolor{block title}{use=structure,fg=structure.fg,bg=structure.fg!20!bg}
	\setbeamercolor{block body}{parent=normal text,use=block title,bg=block title.bg!50!bg, fg=black}
}


% Beginning of a section
\AtBeginSection[]{
	{
		\setbeamercolor{background canvas}{bg=yellow!10}
		\begin{frame}[noframenumbering,plain]
			\framesubtitle{\nameofthepart Chapter \insertromanpartnumber \ -- \iteminsert{\insertpart}}
			\tableofcontents[
				currentsection,
				sectionstyle=show/shaded,
				subsectionstyle=show/show/hide]
		\end{frame}
	\addtocounter{page}{-1}
	%\addtocounter{framenumber}{-1} 
	}
}

% Beginning of a section
\AtBeginSubsection[]{
	{
		\setbeamercolor{background canvas}{bg=green!10}
		\begin{frame}[noframenumbering,plain]
				\framesubtitle{\nameofthepart Chapter \insertromanpartnumber \ -- \iteminsert{\insertpart}}
				\tableofcontents[
					currentsection,
					sectionstyle=show/show,
					currentsubsection,
					subsectionstyle=show/shaded/hide]
			\end{frame}
		\addtocounter{page}{-1}
	}
}


% To have theorems and everything derived prefixed by the slide set
% number
\renewcommand{\thetheorem}{1.\arabic{theorem}}



\title{Matrices are \emph{everywhere}}
\author{Julien Arino}
\date{Fall 2025}

\begin{document}

% The title page
\begin{frame}[noframenumbering,plain]
	\begin{tikzpicture}[remember picture,overlay]
		\node[above right,inner sep=0pt] at (current page.south west)
		{
			\includegraphics[width=\paperwidth]{FIGS/title-page-picture}
		};
	\end{tikzpicture}	
	\titlepage
\end{frame}
\addtocounter{page}{-1}
  
  
\begin{frame}{Outline}
	\tableofcontents[hideallsubsections]
\end{frame}
\addtocounter{page}{-1}


%%%%%%%%%%%%%%%%%%%%%
%%%%%%%%%%%%%%%%%%%%%
%%%%%%%%%%%%%%%%%%%%%
%%%%%%%%%%%%%%%%%%%%%
\begin{frame}
Here are a few examples illustrating how ubiquitous matrices are in mathematics
\vfill
In the process, we introduce some concepts that are used later
\vfill
However, precise definitions are given in subsequent chapter; here concepts are introduced with very few explanations
\end{frame}

%\begin{frame}[noframenumbering,plain]{OUTLINE}
%\tableofcontents
%\end{frame}


%%%%%%%%%%%%%%%%%%%%%
%%%%%%%%%%%%%%%%%%%%%
%%%%%%%%%%%%%%%%%%%%%
%%%%%%%%%%%%%%%%%%%%%
\section{Linear systems of difference equations}
\label{sec:linear_DE}
\begin{frame}{Linear systems of difference equations}
Let $\bx(t)=(x_1(t),\ldots,x_n(t))\in \IR^n$ be a \defword{state variable} and $t\in \IN$ be an \defword{independent variable}, typically thought of as \defword{time}. Let $A\in \mathcal{M}_n(\IR)$. ($\M_n(\IR)$ is also sometimes denoted $\IR^{n\times n}$)
\vfill
An \defword{autonomous homogeneous linear system of difference equations} is a sequence defined by
\begin{subequations}\label{sys:DE}
\begin{align}
\bx(t+1) &= A\bx(t) \label{sys:DE_dx}\\
\bx(0) &=\bx_0\in \IR^n, \label{sys:DE_IC}
\end{align}	
\end{subequations}
where $\bx_0$ is called the \defword{initial condition} (IC)
\vfill
(Autonomous: $A$ is a constant that does not depend on $t$. Homogeneous: the system is not of the form $\bx(t+1)=A\bx(t)+\bb$)
\end{frame}

\begin{frame}
Given the initial condition $\bx(0)=\bx_0$, we have from \eqref{sys:DE_dx} that
\begin{align*}
 \bx(1) &= A\bx(0)=A\bx_0\\
 \bx(2) &= A\bx(1)= AA\bx_0=A^2\bx_0 \\
 \bx(3) &= A\bx(2)=AA^2\bx_0=A^3\bx_0
\end{align*}
\vfill
By induction,
\[
\bx(t)=A^t\bx_0
\] 
for all $t\in\IN$.
Thus the behaviour of $\bx(t)$ as $t\to\infty$ depends on $A^t$
\vfill
In order to understand what this behaviour could be, the following two questions should be answered
\begin{enumerate}
	\item What is (if it exists) $\lim_{t\to \infty} \bx(t)$? 
	\item What is (if it exists) $\lim_{t \to \infty} A^t$? 
\end{enumerate}
\end{frame}

\begin{frame}
Let $\bv$ be an eigenvector associated to the eigenvalue $\lambda$ of $A$, \ie, $\lambda$ be such that $\bv\neq \b0$ satisfies $A\bv=\lambda \bv$. Then we have
 \begin{align*}
 A^2\bv &= A(A\bv) \\
 &=A(\lambda \bv) \\
 &=\lambda A\bv \\
 &=\lambda^2 \bv
 \end{align*}
\ie, $\bv$ is also an eigenvector of $A^2$ and is associated to the eigenvalue $\lambda^2$.
By induction
\[
A^k\bv=\lambda^k \bv
\]
\ie, $\bv$ is an eigenvector of the matrix $A^k$ associated to the eigenvalue $\lambda^k$
\vfill
$A^k\bv$ is a vector in $\IF^n$; we can thus take its norm $\|A^k\bv\|$, where $\|\cdot\|$ is some norm on $\IF^n$. It follows that if $|\lambda|<1$, then $\|A^k\bv\|=|\lambda|^k\|\bv\|$ goes to zero as $k\to\infty$
\end{frame}


\begin{frame}
\begin{theorem}\label{th:TFAE_DE}
Let $A\in\M_n(\IR)$, consider the map $A\bx$. TFAE:
\begin{enumerate}
\item There exists a norm $\|\;\|_\alpha$ on $\IR^n$ and a constant $0<\mu<1$ such that for any $\bx\in\IR^n$, the iterates satisfy, for all $k\geq 0$
\[
\|A^k\bx\|_\alpha\leq \mu^k\|\bx\|_\alpha
\]
\item For any norm $\|\;\|_\beta$ on $\IR^n$, there exists constants $0<\mu<1$ and $C\geq 1$ such that all $\bx\in\IR^n$ and all $k\geq 0$
\[
\|A\bx\|_\beta\leq C\mu^k\|\bx\|_\beta
\]
\item All the eigenvalues $\lambda$ of $A$ satisfy $|\lambda|<1$
\end{enumerate}
\end{theorem}
\vfill
TFAE: ``the following are equivalent''
\end{frame}

\begin{frame}
To use the result above, we need to ensure all eigenvalues lie inside the unit disk in $\IC$
\vfill
For certain classes of matrices, this can be achieved without explicitly computing the eigenvalues
\end{frame}


\begin{frame}{Definitions stemming from Theorem~\ref{th:TFAE_DE}}
A linear map corresponding to a matrix with all eigenvalues of modulus less than 1 is a \defword{linear contraction}, with the origin a \defword{linear sink} or \defword{attracting fixed point}
\vfill
If all eigenvalues have modulus larger than 1, then the map induced by $A$ is a \defword{linear expansion}, and the origin is a \defword{linear source} or \defword{repelling fixed point}
\vfill
The map $A\bx$ is a \defword{hyperbolic linear map} if all eigenvalues of $A$ have modulus different of 1
\end{frame}

\begin{frame}
\begin{definition}
Let $\M_n$ be the set of square $n\times n$ matrices. For $A\in\M_n$, denote $\sigma(A)$ the set of its eigenvalues, \ie,
\[
\sigma(A)=\left\{
\lambda\in\IC; \exists \bv\neq \b0,\; A\bv=\lambda \bv
\right\}
\]
which we call the \defword{spectrum} of $A$. We call \defword{spectral radius} of $A$ the real number
\[
\rho(A)=\max_{\lambda\in\sigma(A)}\{|\lambda|\}
\]
\end{definition}
%%
\vfill
%%
\begin{theorem}\label{th:matrix_everywhere_rho_l_1_goes0}
If $\rho(A)<1$, then $\lim\limits_{t\to \infty} \bx(t)=\b0$ for system \eqref{sys:DE}
\end{theorem}
\end{frame}


%%%%%%%%%%%%%%%%%%%%%%
%%%%%%%%%%%%%%%%%%%%%%
\begin{frame}{Example -- Leslie matrices}
Leslie matrices arise when considering the \emph{age-structured} dynamics of populations reproducing every year, such as most fish
\vfill
Assume $n$ is the maximum age observed for that species. Time $t$ is taken in years and is \emph{discrete}.
Let $\bx(t)=(x_1(t),\ldots,x_n(t))^T$ be the vector of distribution of the population in ages, \ie, $x_i(t)$ is the population of fish of age $a$ between $i-1$ and $i$ at time $t$
\vfill


A proportion $s_i\in[0,1]$ of individuals of age $i-1\leq a<i$ survive to the next year. As years progress at the same rate as age, this means that $x_{i+1}(t+1) = s_i x_i(t)$
\vfill
When individuals reproduce, they give birth to $f_i$ individuals in the first age class, \ie, birth in the first age class takes the form
\[
x_1(t+1)=f_1x_1(t)+f_2x_2(t)+\cdots+f_nx_n(t)
\]
\end{frame}

\begin{frame}{Population growth model with a Leslie matrix}
\begin{align*}
 x_1(t+1) &= f_1x_1(t)+f_2x_2(t)+\cdots+f_nx_n(t) \\
 x_2(t+1) &= s_1x_1(t) \\
 x_3(t+1) &= s_2x_2(t) \\
 &\vdots \\
 x_n(t+1) &= s_{n-1}x_{n-1}(t)
\end{align*}
\end{frame}

\begin{frame}
Suppose that on the first year counts are taken, we observe an initial age distribution $\bx(0)=(x_1(0),\ldots,x_n(0))^T$
\vfill
Assume the survival rates $s_i$ and fecundity constants $f_i$ are known
\vfill
How can we expect the population to evolve over the course of time? In particular, is the species likely to survive or will it become extinct?
\end{frame}

\begin{frame}
Write the system in vector form
\begin{equation}\label{eq:Leslie_long}
\begin{pmatrix}
x_1(t+1) \\ x_2(t+1) \\ x_3(t+1) \\ \vdots \\ x_n(t+1)
\end{pmatrix}
=
\begin{pmatrix}
f_1 & f_2 & f_3 & \cdots & f_{n-1} & f_n \\
s_1 & 0 & 0 & \cdots & 0 & 0 \\
0 & s_2 & 0 & \cdots & 0 & 0 \\
\vdots && \ddots &&& \vdots \\
0 &&&&s_{n-1} & 0
\end{pmatrix}
\begin{pmatrix}
x_1(t) \\ x_2(t) \\ x_3(t) \\ \vdots \\ x_n(t)
\end{pmatrix}
\end{equation}
which we summarize as
\begin{equation}\label{eq:Leslie_short}
\bx(t+1)=L\bx(t)
\end{equation}
with $L$ the matrix in \eqref{eq:Leslie_long}, \ie, a \defword{nonnegative} matrix with only the first row and the first sub-diagonal nonzero
\end{frame}


\begin{frame}
\[
L= \begin{pmatrix}
f_1 & f_2 & f_3 & \cdots & f_{n-1} & f_n \\
s_1 & 0 & 0 & \cdots & 0 & 0 \\
0 & s_2 & 0 & \cdots & 0 & 0 \\
\vdots && \ddots &&& \vdots \\
0 &&&&s_{n-1} & 0
\end{pmatrix}
\]
\vfill
By Theorem~\ref{th:matrix_everywhere_rho_l_1_goes0}, if $\rho(L)<1$ then $\bx(t)\to \b0$ at $t\to\infty$, so in this case, the population would become extinct if $x_i(t)\geq 0$ for all $i,t$
\vfill
By the discussion earlier, if $\rho(L)=1$, then $\|\bx(t)\|$ (the total population) stays constant
\vfill
If $\rho(L)>1$, the population increases
\end{frame}

%%%%%%%%%%%%%%%%%%%%%%
%%%%%%%%%%%%%%%%%%%%%%
%%%%%%%%%%%%%%%%%%%%%%
%%%%%%%%%%%%%%%%%%%%%%
\section{Ordinary differential equation}

\begin{frame}{ODE}
An \defword{ordinary differential equation} (ODE) is an equation of the form 
\begin{equation}\label{eq:nonlinear_ODE}
\dfrac{d}{dt}\bx(t)= f(\bx(t))
\end{equation}
where $\bx(t)\in \IR^n$ is a \emph{function} and $f:\IR^n \to \IR^n$ is a \defword{vector field}
\vfill
An \defword{initial value problem} (IVP) is the consideration of \eqref{eq:nonlinear_ODE} together with an initial condition $\bx(t_0)=\bx_0\in\IR^n$
\vfill
A solution to \eqref{eq:nonlinear_ODE} is a function $\phi(t)$ that satisfies \eqref{eq:nonlinear_ODE}. A solution to the IVP with $\bx(t_0)=\bx_0$ associated to \eqref{eq:nonlinear_ODE} is, among all solutions to \eqref{eq:nonlinear_ODE}, the one (typically the only one) that additionally satisfies the initial condition, \ie, such that $\phi(t_0)=\bx_0$
\end{frame}

\begin{frame}{Linear systems}
Consider the autonomous linear system
\begin{equation}\label{eq:Lin_EDO_cc_hom}
\bx'(t)=A\bx(t)
\end{equation}
where $A\in\M_n$ is a constant
\vfill
Let $\sigma(A)=\{\lambda_1,\ldots,\lambda_n\}$ be the spectrum of $A$.
Let $\bw_j=\bu_j+i\bv_j$ be a generalized eigenvector of $A$ corresponding to an eigenvalue $\lambda_j=a_j+ib_j$, with $\bv_j=\b0$ if $b_j=0$, and 
\[
B=\{\bu_1,\ldots,\bu_k,\bu_{k+1},\bv_{k+1},\ldots,\bu_m,\bv_m\}
\]
be a basis of $\mathbb{R}^n$, with $n=2m-k$, where $k$ is the number of real eigenvalues in $\sigma(A)$
\end{frame}

\begin{frame}
\begin{definition}[Stable, unstable and center subspaces]
The \defword{stable}, \defword{unstable} and \defword{center} subspaces of the linear system \eqref{eq:Lin_EDO_cc_hom} are given, respectively, by
\[
E^s=\textrm{Span}\{u_j,v_j:a_j<0\},
\]
\[
E^u=\textrm{Span}\{u_j,v_j:a_j>0\}
\]
and
\[
E^c=\textrm{Span}\{u_j,v_j:a_j=0\}
\]
\end{definition}
\end{frame}

\begin{frame}
\begin{definition}
The mapping $e^{At}:\mathbb{R}^n\to\mathbb{R}^n$ is the \defword{flow} of the linear system \eqref{eq:Lin_EDO_cc_hom}
\end{definition}
The term \defword{flow} is used since $e^{At}$ describes the motion of points $\bx_0\in\mathbb{R}^n$ along trajectories of \eqref{eq:Lin_EDO_cc_hom}
%%
\vfill
%%
\begin{definition}
If all eigenvalues of $A$ have nonzero real part, \ie, if $E^c=\emptyset$, then the flow $e^{At}$ of system \eqref{eq:Lin_EDO_cc_hom} is a \defword{hyperbolic flow} and the system \eqref{eq:Lin_EDO_cc_hom} is a \defword{hyperbolic linear system}
\end{definition}
%%
\vfill
%%
\begin{definition}
A subspace $E\subset\mathbb{R}^n$ is \defword{invariant with respect to the flow} $e^{At}$, or \defword{invariant under the flow} of \eqref{eq:Lin_EDO_cc_hom}, if $e^{At}E\subset E$ for all $t\in\mathbb{R}$
\end{definition}
\end{frame}

\begin{frame}
\begin{theorem}
Let $E$ be the generalized eigenspace of $A$ associated to the eigenvalue $\lambda$. Then $AE\subset E$
\end{theorem}
%%
\vfill
%%
\begin{theorem}
Let $A\in\mathcal{M}_n(\mathbb{R})$. Then
\[
\mathbb{R}^n=E^s\oplus E^u\oplus E^c
\]
Furthermore, if the matrix $A$ is the matrix of the linear autonomous system \eqref{eq:Lin_EDO_cc_hom}, then $E^s$, $E^u$ and $E^c$ are invariant under the flow of \eqref{eq:Lin_EDO_cc_hom}, \ie, let
$x_0\in E^S$, $y_0\in E^C$ and $z_0\in E^U$, then $e^{At}x_0\in E^S$,
$e^{At}y_0\in E^C$ and $e^{At}z_0\in E^U$
\end{theorem}
\end{frame}



\begin{frame}{Nonlinear systems of ODE}
There is no general theory allowing to obtain explicit solutions of a nonlinear IVP
\vfill
Instead of seeking explicit solutions, we use \defword{qualitative analysis}, which uses analysis to establish properties of the solutions without needing to actually find their explicit expression
\end{frame}


\begin{frame}
Suppose you are given a system of ordinary differential equations
\begin{equation}\label{sys:ODE}
\bx'=f(\bx),
\end{equation}
where $\bx\in\IR^n$ and $f:\IR^n\to\IR^n$ is $C^1$. A standard step when studying \eqref{sys:ODE} \defword{qualitatively} is to seek \defword{equilibria} of \eqref{sys:ODE}, \ie, points $\bx^\star\in\IR^n$ such that
\begin{equation}\label{eq:equilibrium}
f(\bx^\star)=\b0.
\end{equation}
At such a point, $\bx'=\b0$, meaning that system \eqref{sys:ODE} is at rest. If you were to consider solutions to \eqref{sys:ODE} with an initial condition $\bx(0)=\bx^\star$, then there would hold that $\bx(t)=\bx^\star$ for all $t\geq 0$
\vfill
What would happen if instead of starting at $\bx^\star$, you were to choose an initial condition $\bx(0)$ close to but distinct from $\bx^\star$?
\end{frame}

\begin{frame}{Simplified version of Hartman-Grobman}
\begin{definition}
An equilibrium point $\bx^\star$ is \defword{hyperbolic} if the Jacobian matrix $Df$ of \eqref{sys:ODE} evaluated at $\bx^\star$, denoted $Df(\bx^\star)$, has no eigenvalues with zero real part, \ie, $Df(\bx^\star)$ is invertible
\end{definition}
\vfill
\begin{theorem}[Hartman-Grobman]\label{th:HartmanGrobman}
Let $\bx^\star$ be a hyperbolic equilibrium point of \eqref{sys:ODE}.
Then in some neighbourhood $\mathcal{N}(\bx^\star)$ of $\bx^\star$, the flow of \eqref{sys:ODE} is topologically equivalent to the flow of the linear system 
\begin{equation}\label{sys:linearized}
\bx'=Df(\bx^\star)(\bx-\bx^\star)
\end{equation} 
where $Df(\bx^\star)$ is the Jacobian matrix $Df$ of $f$ evaluated at $\bx^\star$
\end{theorem}
\end{frame}

\begin{frame}
\begin{theorem}[Stable manifold theorem]\label{th:stable_manifold}
Let $E$ be an open subset of $\IR^n$ containing the origin, let $f\in C^1(E)$ and let $\phi_t$ be the flow of the nonlinear system \eqref{sys:ODE}. 
Suppose that $f(\b0)=\b0$ and that $Df(\b0)$ has $k$
eigenvalues with negative real part and $n-k$ eigenvalues with positive real part. 
Then there exists a $k$-dimensional differentiable
manifold $S$ tangent to the stable subspace $E^s$ of the linear system \eqref{sys:linearized} at $\b0$ such that for all $t\geq 0$, $\phi_t(S)\subset S$ and for all $\bx_0\in S$
\[
\lim_{t\to\infty} \phi_t(\bx_0)=\b0
\]
and there exists an $(n-k)$-dimensional differentiable manifold $U$ tangent to the unstable subspace $E^u$ of (\ref{sys:linearized}) at $\b0$ such that for all $t\leq 0$, $\phi_t(U)\subset U$ and for all $\bx_0\in U$
\[
\lim_{t\to-\infty}\phi_t(\bx_0)=\b0
\]
\end{theorem}
%%
\end{frame}
%%%%%%%%%%%%%%%%
%%%%%%%%%%%%%%%%

%%%%%%%%%%%%%%%%
%%%%%%%%%%%%%%%%
\begin{frame}{Example -- A chemostat model}
System of 2 nonlinear DE modelling a biological device called a \defword{chemostat}
\begin{subequations}
\label{sys:chemostat}
\begin{align}
\frac{dS}{dt} &= D\left(S^0-S\right)-\mu(S)x \label{sys:chemostatS} \\
\frac{dx}{dt} &= (\mu(S)-D)x \label{sys:chemostatx}
\end{align}
\end{subequations}
%%
\vfill
%%
Parameters $S^0$ and $D$, respectively the \defword{input concentration} and the \defword{dilution} rate, are real and positive.
The function $\mu$ is the \defword{growth} function. It is generally
assumed to satisfy $\mu(0)=0$, $\mu'>0$ and $\mu''<0$
\end{frame}

\begin{frame}
One can verify that the positive quadrant is
positively invariant under the flow of (\ref{sys:chemostat}),
\ie, that for $S(0)\geq 0$ and $x(0)\geq 0$, solutions remain nonnegative for all positive times, and similar properties
\vfill
But since we are here only interested in applications of the stable
manifold theorem, we proceed to a very crude analysis, and will not
deal with this point
\end{frame}

\begin{frame}
Write the system in vector form as
\[
\xi'=f(\xi),
\]
with $\xi=(S,x)^T$ and
\[
f(\xi)=
\begin{pmatrix}
D(S^0-S)-\mu(S)x \\
(\mu(S)-D)x
\end{pmatrix}
\]
EP of the system are found by solving $f(\xi)=0$. We find one situated on one of the boundaries of the positive quadrant
\[
\xi_T^\star=(S^\star_T,x^\star_T)=\left(S^0,0\right)
\]
and the second one in the interior of $\IR_+^2$,
\[
\xi_I^\star=(S^\star,x^\star)=\left(\lambda,S^0-\lambda\right)
\]
where $\lambda$ is such that $\mu(\lambda)=D$. Note that this implies
that if $\lambda\geq S^0$, $\xi_T^\star$ is the only equilibrium of the
system since in that case, $\xi_I^\star\not\geq 0$, which is not biologically realistic
\end{frame}

\begin{frame}
At an arbitrary point $\xi=(S,x)$, the Jacobian matrix of \eqref{sys:chemostat} is given by
\[
Df(\xi)=\left(
\begin{matrix}
-D-\mu'(S)x & -\mu(S) \\
\mu'(S)x & \mu(S)-D
\end{matrix}
\right)
\]
%%
\vfill
%%
Thus, at the trivial equilibrium $\xi_T^\star$
\[
Df(\xi_T^\star)=
\left(
\begin{matrix}
-D & -\mu(S^0) \\
0 & \mu(S^0)-D
\end{matrix}
\right)
\]
%%
\vfill
%%
We have two eigenvalues, $-D$ and $\mu(S^0)-D$.
Since $-D<0$, we focus on the eigenvalue $\mu(S^0)-D$
\end{frame}

\begin{frame}
Assume that
$\mu(S^0)-D<0$. This implies that $\xi_T^\star$ is the only
equilibrium, since $\xi_I^\star$ is not feasible
if $\lambda>S^0$
\vfill
System has dimensionality 2 and $Df(\xi_T^\star)$
has two negative eigenvalues $\implies$ the stable manifold theorem
(Theorem~\ref{th:stable_manifold}) states that there exists a
2-dimensional differentiable manifold $\mathcal{M}$ such that
\begin{itemize}
\item $\phi_t(\mathcal{M})\subset\mathcal{M}$,
\item for all $\xi_0\in\mathcal{M}$,
$\lim_{t\to\infty}\phi_t(\xi_0)=\xi_T^\star$.
\item At $\xi_T^\star$, $\mathcal{M}$ is tangent to the stable subspace
$E^S$ of the linearized system $\xi'=Df(\xi_T^\star)(\xi-\xi_T^\star)$.
\end{itemize}
Since there are no eigenvalues with positive real part, there does not
exist an unstable manifold in this case 
\end{frame}

\begin{frame}
Let us now characterize the
nature of the stable subspace $E^S$. It is obtained by studying the
linear system
\begin{align}
\xi' &= Df(\xi_T^\star)(\xi-\xi_T^\star) \nonumber \\
&= \left(
\begin{matrix}
-D & -\mu(S^0) \\
0 & \mu(S^0)-D
\end{matrix}
\right)
\left(
\begin{matrix}
S-S^0 \\
x
\end{matrix}
\right) \nonumber \\
&= \left(
\begin{matrix}
-D(S-S^0)-\mu(S^0)x \\
(\mu(S^0)-D)x
\end{matrix}
\right)
\end{align}
Of course, the Jacobian matrix associated to this system is the same
as that of the nonlinear system (at $\xi_T^\star$). Associated to the
eigenvalue $-D$ is the eigenvector $v_1=(1,0)^T$, to $\mu(S^0)-D$ is
$v_2=(-1,1)^T$
\end{frame}

\begin{frame} 
The stable subspace is thus given by $\Span(v_1,v_2)$, \ie,
the whole of $\IR^2$
\vfill
In fact, the stable manifold of $\xi_T^\star$ is the whole positive
quadrant, since all solutions limit to this equilibrium 
\vfill
The same type of analysis can be conducted at the interior equilibrium $\xi_I^\star$. It is a little harder in this case, since $x^\star>0$ there and therefore the Jacobian matrix $Df(\xi_I^\star)$ does not have the same upper triangular structure as $Df(\xi_T^\star)$
\end{frame}

%%%%%%%%%%%%%%
%%%%%%%%%%%%%%
%%%%%%%%%%%%%%
%%%%%%%%%%%%%%
\section{Linear systems of ODE -- Brief theory}


\frame{\frametitle{Linear ODEs}
\begin{definition}[Linear ODE]
A \emph{linear} ODE is a differential equation taking the form
\begin{equation}\label{sys:linear_general}
\frac{d}{dt}\bx=A(t)\bx+\bB(t),\tag{LNH}
\end{equation}
where $A(t)\in\mathcal{M}_n(\IR)$ with continuous entries, $\bB(t)\in\IR^n$ with real valued, continuous coefficients, and $\bx\in\IR^n$
\vskip0.5cm
The associated IVP takes the form 
\begin{equation}\label{sys:IVP_linear_general}
\begin{aligned}
\frac{d}{dt}\bx &= A(t)\bx+\bB(t) \\
\bx(t_0)&=\bx_0
\end{aligned}
\end{equation}
\end{definition}
}

\frame{\frametitle{Types of linear systems}
\begin{itemize}
\item $\bx'=A(t)\bx+\bB(t)$ is linear nonautonomous ($A(t)$ depends on $t$) nonhomogeneous (also called \emph{affine} system)
\vfill
\item $\bx'=A(t)\bx$ is linear nonautonomous homogeneous
\vfill
\item $\bx'=A\bx+\bB$, that is, $A(t)\equiv A$ and $\bB(t)\equiv \bB$, is linear autonomous nonhomogeneous (or affine autonomous)
\vfill
\item $\bx'=A\bx$ is linear autonomous homogeneous
\end{itemize}
}


\frame{\frametitle{Existence and uniqueness of solutions}
\begin{theorem}[Existence and Uniqueness]
Solutions to \eqref{sys:IVP_linear_general} exist and are unique on the whole interval over which $A$ and $\bB$ are continuous
\vskip0.5cm
In particular, if $A,\bB$ are constant, then solutions exist on $\IR$
\end{theorem}
}

\frame{\frametitle{Autonomous linear systems}
Consider the autonomous affine system
\begin{equation}\label{sys:affine_auton}
\frac{d}{dt}\bx=A\bx+\bB\tag{A}
\end{equation}
and the associated homogeneous autonomous system
\begin{equation}\label{sys:lin_auton}
\frac{d}{dt}\bx=A\bx\tag{L}
\end{equation}
}

\frame{\frametitle{Exponential of a matrix}
\begin{definition}[Matrix exponential]
Let $A\in\M_n(\IF)$ with $\IF=\IR$ or $\IC$. 
The \defword{matrix exponential} of $A$, denoted $e^{At}$, is a matrix in $\M_n(\IF)$, defined by
\[
e^{At}=\II +\sum_{k=1}^\infty \frac{t^k}{k!}A^k
\]
where $\II$ is the identity matrix in $\M_n(\IF)$
\end{definition}
}

\begin{frame}\frametitle{Properties of the matrix exponential}
\begin{itemize}
\item $e^{At_1}e^{At_2}=e^{A(t_1+t_2)}$ for all $t_1,t_2\in\IR$
\vfill
\item $Ae^{At}=e^{At}A$ for all $t\in\IR$
\vfill
\item $(e^{At})^{-1}=e^{-At}$ for all $t\in\IR$
\vfill
\item The unique solution $\phi$ of \eqref{sys:lin_auton} with $\phi(t_0)=\bx_0$ is given by
\[
\phi(t)=e^{A(t-t_0)}\bx_0
\]
\end{itemize}
\end{frame}

\begin{frame}\frametitle{Using the matrix exponential}
Let $P$ be a nonsingular matrix in $\M_n(\IR)$. We transform the IVP
\begin{equation}\label{IVP:lin_auton}
\begin{aligned}
\frac{d}{dt}\bx &= A\bx\\
\bx(t_0)&=\bx_0
\end{aligned}\tag{LIVP}
\end{equation}
using the transformation $\bx=P\by$ or $\by=P^{-1}\bx$
\end{frame}

\begin{frame}
The dynamics of $\by$ is $\by'=P^{-1}AP\by$
\begin{align*}
\by' &= (P^{-1}\bx)' \\
&= P^{-1}\bx' \\
&= P^{-1}A\bx \\
&= P^{-1}AP\by
\end{align*}
\vfill
The initial condition is $\by_0=P^{-1}\bx_0$
\end{frame}

\begin{frame}
We have thus transformed IVP \eqref{IVP:lin_auton} into
\begin{equation}\label{IVP:lin_auton2}
\begin{aligned}
\frac{d}{dt}\by &= P^{-1}AP\by\\
\by(t_0)&=P^{-1}\bx_0
\end{aligned}\tag{LIVP$_y$}
\end{equation}
\vfill
From the earlier result, we then know that the solution of \eqref{IVP:lin_auton2} is given by
\[
\psi(t)=e^{P^{-1}AP(t-t_0)}P^{-1}\bx_0
\]
and since $\bx=P\by$, the solution to \eqref{IVP:lin_auton} is given by
\[
\phi(t)=Pe^{P^{-1}AP(t-t_0)}P^{-1}\bx_0
\]
\vfill
So everything depends on $P^{-1}AP$
\end{frame}

\frame{\frametitle{The cases}
\begin{itemize}
\item $P^{-1}AP$ is diagonal,
the solution to \eqref{IVP:lin_auton} is given by
\[
\phi(t)=P
\begin{pmatrix}
e^{\lambda_1t} & & 0 \\
& \ddots & \\
0 & & e^{\lambda_nt}
\end{pmatrix}
P^{-1}x_0
\]
\vfill
\item $P^{-1}AP$ is not diagonal, then use Jordan form (slightly more complicated)
\end{itemize}
}

\frame{
\begin{theorem}
For all $(t_0,\bx_0)\in\IR\times\IR^n$, there is a unique solution $\bx(t)$ to \eqref{IVP:lin_auton} defined for all $t\in\IR$. Each coordinate function of $\bx(t)$ is a linear combination of functions of the form
\[
t^ke^{\alpha t}\cos(\beta t)\quad\textrm{and}\quad t^ke^{\alpha t}\sin(\beta t)
\]
where $\alpha+i\beta$ is an eigenvalue of $A$ and $k$ is less than the algebraic multiplicity of the eigenvalue
\end{theorem}
}

\frame{\frametitle{Generalized eigenvectors, nilpotent matrix}
\begin{definition}[Generalized eigenvectors]
Let $A\in\M_r(\IR)$. Suppose $\lambda$ is an eigenvalue of $A$ with multiplicity $m\leq n$. Then, for $k=1,\ldots,m$, any nonzero solution $\bv$ of
\[
(A-\lambda\II)^k\bv=\b0
\]
is called a \emph{generalized eigenvector} of $A$
\end{definition}
\vfill
\begin{definition}[Nilpotent matrix]
Let $A\in\M_n(\IR)$. $A$ is \emph{nilpotent} (of order $k$) if $A^{j}\neq 0$ for $j=1,\ldots,k-1$, and $A^k=\b0$
\end{definition}
}

\frame{\frametitle{Jordan normal form}
\begin{theorem}[Jordan normal form]
Let $A\in\M_n(\IR)$ have eigenvalues $\lambda_1,\ldots,\lambda_n$, repeated according to their multiplicities
\begin{itemize}
\item 
Then there exists a basis of generalized eigenvectors for $\IR^n$
\item 
And if $\{v_1,\ldots,v_n\}$ is any basis of generalized eigenvectors for $\IR^n$, then the matrix $P=[v_1\cdots v_n]$ is invertible, and $A$ can be written as
\[
A=S+N
\]
where
\[
P^{-1}SP=\diag(\lambda_j)
\]
the matrix $N=A-S$ is nilpotent of order $k\leq n$ and $S$ and $N$ commute, i.e., $SN=NS$
\end{itemize}
\end{theorem}
}

\frame{
\begin{theorem}
Under conditions of the Jordan normal form Theorem, the linear system $\bx'=A\bx$ with initial condition $\bx(0)=\bx_0$, has solution
\[
\bx(t)=P\diag\left(e^{\lambda_j t}\right)P^{-1}\left(\II+Nt+\cdots\frac{t^k}{k!}N^k\right) \bx_0
\]
\end{theorem}
\vfill
\begin{theorem}[Case of an eigenvalue of multiplicity $n$]
Suppose that $\lambda$ is an eigenvalue of multiplicity $n$ of $A\in\M_n(\IR)$. Then $S=\diag(\lambda)$, and the solution of $\bx'=A\bx$ with initial value $\bx_0$ is given by
\[
\bx(t)=e^{\lambda t}\left(\II+Nt+\cdots\frac{t^k}{k!}N^k\right) \bx_0
\]
\end{theorem}
In the simplified case, we do not need the basis of generalized eigenvectors (i.e., $P$)
}

 \frame{\frametitle{A variation of constants formula}
 \begin{theorem}[Variation of constants formula]
 Consider the IVP
 \begin{subequations}\label{sys:lin_nonlin}
 \begin{align}
 x' &= Ax+B(t) \\
 \bx(t_0) &= \bx_0
 \end{align}
 \end{subequations}
 where $B:\IR\to\IR^n$ a smooth function on $\IR$ and $e^{A(t-t_0)}$ is the matrix exponential associated to the homogeneous system $\bx'=A\bx$
 \vskip0.5cm
 Then the solution $\phi$ to \eqref{sys:lin_nonlin} is given by
 \begin{equation}
 \phi(t)=e^{A(t-t_0)}\bx_0 + \int_{t_0}^t e^{A(t-s)}B(s)ds
 \end{equation}
 \end{theorem}
 }

%%%%%%%%%%%%%%%%%%%%%%
%%%%%%%%%%%%%%%%%%%%%%
%%%%%%%%%%%%%%%%%%%%%%
%%%%%%%%%%%%%%%%%%%%%%
\section{Discrete-time Markov chains}
\begin{frame}{Discrete-time Markov chain}
A \defword{discrete-time Markov chain} is a stochastic process
\vfill
Consider a system with $n$ \defword{states} denoted $S_1,\ldots,S_n$. The system starts in a given state. Every time step, it switches to a different state. (Transition from one state to itself is also allowed)
\vfill
We assume that the system is \emph{stochastic}, \ie, that the transitions happen at random. In discrete-time Markov chains, the instants at which transitions occur (or not) are in a discrete set, typically rescaled to be $\IN$
\vfill
(In \emph{continuous-time} Markov chains, the time of the switches themselves is a continuous random variable, so times are in $\IR$)
\end{frame}

\begin{frame} 
Let $p_i(t)$ be the probability that state $S_i$ occurs on the $t^{th}$ time step, $1\leq i\leq n$
\vfill
One of the key assumptions in Markov chains is that the process is \defword{memoryless}: the transition that occurs from time $t$ to time $t+1$ depends only on the state of the system at time $t$
\vfill
Since one the states $S_i$ must occur on the $t^{th}$ time step,
\[
p_1(t)+p_2(t)+\cdots+p_n(t)=1
\]
\end{frame}

\begin{frame}{The transition matrix}
Let
\[
P=[p_{ij}]
\]
be the \defword{transition matrix} (or \defword{projection matrix}) of the Markov chain, where
\[
p_{ij} = \IP(S_i|S_j)
\]
\ie, the probability of making a transition from state $j$ to state $i$
\vfill
Many texts, \emph{e.g.}, \cite{KemenySnell1983}, define $p_{ij}=\IP(S_j|S_i)$. This works the same way, except it leads to $P^T$ instead of $P$ (or row vectors instead of column vectors)
\end{frame}

\begin{frame} 
Let $p_i(t+1)$ be the probability that state $S_i$, $1\leq i\leq n$, occurs on the $(t+1)^{th}$ step. 
There are $n$ ways to be in state $S_i$ at step $t+1$:
\begin{enumerate}
	\item Step $t$ is $S_1$. The probability of getting $S_1$ on the $t^{th}$ step is $p_1(t)$ and the probability of having $S_i$ after $S_1$ is $p_{i1}$, so $\IP(S_i|S_1)=p_{i1}p_1(k)$
	\item We get $S_2$ on step $t$ and $S_i$ on step $t+1$ and $\IP(S_i|S_2)=p_{i2}p_2(t)$
	\item[..]
	\item[n.] The probability of occurrence of $S_i$ at step $t+1$ if $S_n$ at step $t$ is $\IP(S_i|S_n)=p_{in}p_n(t)$
\end{enumerate}
\vfill
\begin{align*}
\implies p_i(t+1) &= \IP(S_i|S_1)+\cdots+\IP(S_i|S_n) \\
&= p_{i1}p_1(t)+\cdots+p_{in}p_n(t)
\end{align*}
\end{frame}

\begin{frame} 
Therefore,
\begin{align*}
p_1(t+1) &= p_{11}p_1(t)+p_{12}p_2(t)+\dots+p_{1n}p_n(t) \\
& \;\;\vdots\\
p_k(t+1) &= p_{k1}p_1(t)+p_{k2}p_2(t)+\dots+p_{kr}p_r(t) \\
& \;\;\vdots\\
p_n(t+1) &= p_{n1}p_1(t)+p_{n2}p_2(t)+\dots+p_{nn}p_r(t)
\end{align*}
\end{frame}

\begin{frame} 
In matrix form,
\begin{equation}\label{eq:CTMC_vector}
p(t+1)=Pp(t), \quad t=1,2,3,\dots
\end{equation}
where $p(t)=(p_1(t),p_{2}(t),\dots , p_n(t))^T$ is a (column) probability vector and $P=[p_{ij}]\in\M_n$ is the transition matrix
\begin{equation}\label{eq:DMTC_transition_matrix}
P=
\begin{pmatrix}
p_{11} & p_{12} & \cdots & p_{1n} \\
p_{21} & p_{22} & \cdots & p_{2n} \\
\vdots & \vdots && \vdots \\
p_{n1} & p_{n2} & \cdots & p_{nn}
\end{pmatrix}
\end{equation}
\end{frame}


\begin{frame}{Two observations about $P$}
\begin{itemize}
	\item Entries of $P$ being probabilities, they are all in $[0,1]$. So in particular, they are all nonnegative. We say $P$ is a \defword{nonnegative matrix} and write $P\geq 0$
	\vfill
	\item The \emph{column} sums of $P$ all equal 1. Take for instance the first column: its entries represent the probabilities of transition from state $1$. Since an event must always happen, the sum of these probabilities \emph{must} be 1. We say $P$ is a (column) \defword{stochastic matrix}
\end{itemize}
\end{frame}

\begin{frame} 
\begin{definition}[Stochastic matrix]
	The nonnegative $n\times n$ matrix $M$ is
	\begin{itemize}
		\item \defword{row stochastic} if $\sum_{j=1}^nm_{ij}=1$ for all $i=1,\dots, n$
		\item \defword{column stochastic} if $\sum_{i=1}^nm_{ij}=1$ for all $j=1,\dots, n$
		\item \defword{stochastic} if it is row or column stochastic
		\item \defword{doubly stochastic} if it is row \emph{and} column stochastic
	\end{itemize}
\end{definition}
\end{frame}


\begin{frame}{An interesting property of stochastic matrices}
\begin{theorem}\label{th:spectrum_stochastic_matrix}
	Let $M$ be a stochastic matrix. Then all eigenvalues $\lambda$ of $M$ are such that $|\lambda|\leq 1$. 
	Furthermore, $\lambda =1$ is an eigenvalue of $M$
\end{theorem}
%%
\vfill
%%
Assume w.l.o.g. that $M$ is row stochastic, \ie, $M$ has row sums 1. In vector form, $M\nbOne=\nbOne$. Now remember that $\lambda$ is an eigenvalue of $M$, with associated eigenvector $\bv\neq \b0$, if and only if $M\bv=\lambda \bv$. So, in the expression $M\nbOne=\nbOne$, we read an eigenvector, $\nbOne$, and an eigenvalue, $1$
\vfill
Proving the first conclusion of Theorem~\ref{th:spectrum_stochastic_matrix} involves a theorem called the Perron-Frobenius Theorem, which we will see in much detail later
\end{frame}

\begin{frame}{Long time behaviour}
Let $p_0:=p(0)=(p_1(0),\ldots,p_n(0))^T$ be the initial probability distribution vector, with $\nbOne^Tp_0=1$, \ie, such that the sum of the entries of $p_0$ be 1; we could also write $\langle p_0,\nbOne\rangle = 1$. Then
\begin{align*}
p(1) &= Pp(0) = Pp_0 \\
p(2) &= Pp(1)\\
&= P(Pp_0) \\
&= P ^2p_0
\end{align*}
Iterating, for any $t\in\IN$,
\[
p(t)=P^tp_0
\]
\vfill
Rings a bell? (cf Section~\ref{sec:linear_DE})
\end{frame}

\begin{frame}
So the long time evolution of the system is governed by
\begin{equation}\label{eq:limit_DTMC}
\lim_{t\rightarrow +\infty}p(t)
=\lim_{t\rightarrow +\infty}P^tp_0
=\left(\lim_{t\rightarrow +\infty}P^t\right)p_0
\end{equation}
if the latter limit exists
\vfill
So if we can characterize the nature of matrix $P^t$ and in particular, the existence of the limit $\lim_{t\to\infty}P^t$, we will know the long time behaviour of the Markov chain
\end{frame}

\begin{frame}
\begin{theorem}
If $M,N$ are nonsingular stochastic matrices, then $MN$ is a stochastic matrix
\end{theorem}
\vfill
\begin{corollary}
If $M$ is a nonsingular stochastic matrix, then for any $t\in\IN$, $M^t$ is a stochastic matrix
\end{corollary}
\vfill
$\implies$ matrix $P^t$ in \eqref{eq:limit_DTMC} is stochastic; so, in particular, it is a nonnegative matrix with column sums all equal to 1
\end{frame}


\begin{frame}{Regular Markov chains}
\begin{definition}[Regular Markov chain]
A \defword{regular Markov chain} is one in which $P^k$ is positive for some integer $k>0$, i.e., $P^k$ has only positive entries, no zero entries
\end{definition}
\vfill
\begin{definition}[Primitive matrix]
A nonnegative matrix $M$ is \defword{primitive} if, and only if, there is an integer $k>0$ such that $M^k$ is (entry-wise) positive
\end{definition}
\vfill
\begin{theorem}
A Markov chain is regular $\iff$ $P$ is primitive
\end{theorem}
\end{frame}

\begin{frame}{Regular Markov chains are well-behaved}
\begin{theorem}\label{th:behaviour_regular_MC}
If $P$ is the transition matrix of a regular Markov chain, then
\begin{enumerate}
\item the powers $P^t$ approach a stochastic matrix $W$
\item each column of $W$ is the same vector $w=(w_1,\ldots,w_n)^T$
\item the components of $w$ are positive
\end{enumerate}
\end{theorem}
\vfill
So if the Markov chain is regular, \eqref{eq:limit_DTMC} becomes
\[
\lim_{t\rightarrow +\infty}p(t)=\lim_{t\rightarrow +\infty}P^tp_0
=Wp_0
\]
\end{frame}

\begin{frame} 
Let $M\in\M_n$, $u,v$ be two column vectors, $\lambda\in\IC$. Then, if  
\[
Mu=\lambda u
\]
$u$ is the (right) eigenvector corresponding to $\lambda$, and if
\[
\bv^TM=\lambda \bv^T
\]
then $\bv$ is the left eigenvector corresponding to $\lambda$
\vfill 
To a given eigenvalue there corresponds one left and one right eigenvector (to multiples)
\vfill
$(\bv^TM)=(\lambda \bv^T)^T\iff M^T\bv=\lambda \bv$, so if your numeric/symbolic solver spits out right eigenvectors, to get left ones, compute eigenvectors of $M^T$
\end{frame}

\begin{frame}{Back to the regular MC}
We already know that the left eigenvector corresponding to 1 is $\nbOne^T$, since $\nbOne^TP=\nbOne^T$, \ie, the column sums of $P$ all equal 1
\vfill
The vector $w$ in Theorem~\ref{th:behaviour_regular_MC} is the right eigenvector corresponding to the eigenvalue 1 of $P$
\vfill
To see this, remark that, if $p(t)$ converges, then $p(t+1)=Pp(t)$ in the limit for large $t$, so $w$ is a fixed point of the system. We thus write
\[
w=Pw
\]
and solve for $w$, which amounts to finding $w$ as the (right) eigenvector corresponding to the eigenvalue 1
\end{frame}

\begin{frame}
When you compute an eigenvector, the result is to a multiple and often the expression needs normalising (you want a probability vector)
\vfill
Once you obtain $w$, check that the norm 
\[
\|w\|=w_1+\cdots+w_n
\] 
is equal to one. If not, normalise as 
\[
\frac{w}{\|w\|}
\]
\end{frame}


\begin{frame}{Absorbing Markov chains}
Suppose now that the matrix is not only not primitive but also \defword{reducible}
\vfill
\begin{definition}[Reducible/irreducible matrices]
$0\leq M\in\M_n$ is \defword{reducible} if $\exists P\in\M_n$, permutation matrix, s.t.
\[
P^TMP =
\begin{pmatrix}
S & R \\ 0 & Q
\end{pmatrix}
\]
If no such matrix exists, $M$ is \defword{irreducible}
\end{definition}
\vfill
\begin{theorem}
Let $\G(M)$ be the digraph induced by $M$. We have $0\leq M\in\M_n$ irreducible $\iff$ $\G(M)$ strongly connected
\end{theorem}
\end{frame}

\begin{frame}
So if the transition matrix $P$ is reducible, $\G(P)$ is not strongly connected, \ie, there are states of the chain that are not accessible from others
\vfill
When in a state that does not have access to other states, we are stuck.. we say \defword{absorbed}
\vfill
An \defword{absorbing Markov chain} is one where at least one state is absorbing
\vfill
Theorem~\ref{th:behaviour_regular_MC} does not apply here, but we get a lot of interesting properties by observing that $P$ can be put in the form
\[
\begin{pmatrix}
\II& R \\ 0 & Q
\end{pmatrix}
\]
and considering the \defword{fundamental matrix} $N=(\II-Q)^{-1}$
\end{frame}

%%%%%%%%%%%%%%%%%%
%%%%%%%%%%%%%%%%%%
%%%%%%%%%%%%%%%%%%
%%%%%%%%%%%%%%%%%%
%%%%%%%%%%%%%%%%%%%%%%%%%%%%%%%%
%%%%%%%%%%%%%%%%%%%%%%%%%%%%%%%%
\section{Discretisation of partial differential equations}
\label{sec:discretisation_PDEs}

\begin{frame}
Following \cite{Windisch1989}: $\Omega\subset\IR^d$ ($d\geq 1$) a bounded connected region, $u=u(x)$ for $x\in\Omega$; a linear elliptic BVP on $\Omega$ takes the form 
\begin{subequations}\label{sys:PDE}
\begin{align}
Lu &= f,\quad \Omega \\
\alpha u+\beta \frac{\partial u}{\partial v} &= g, \quad \partial\Omega
\end{align}
\end{subequations}
where $L$ is a linear differential operator of the form
\[
Lu = -\nabla(k\nabla u+bu)+qu
\]
Suppose that the scalar functions $k(x)$, $q(x)$ and the vector function
$b(x)=(b_1(x),\ldots,b_d(x))^T$ for $d>1$, otherwise a scalar function
$b(x)$, are sufficiently smooth over the region $\Omega$. 
Further, let $k(x)\geq k_0=\textrm{const} > 0$ and $q(x)\geq 0$ for each $x\in\Omega$
\vfill
$L$ together with corresponding BC on $\partial\Omega$ is a linear elliptic differential operator
\end{frame}


\begin{frame} 
As a particular example, suppose $\Omega=(0,1)$ is a one-dimensional domain, i.e., an interval, and suppose
\[
\bar{\omega}_h=
\left\{
0=x_1<x_2<\ldots<x_n=1
\right\}
=\omega_h+\gamma_h
\]
is a discretisation of the closure of said interval, with $\gamma_h=\{x_0,x_n\}$ and $h_i=h=1/n$ be a uniform grid, i.e., $x_i=ih$. Consider the following problem
\begin{subequations}
\begin{align}
Lu &= -u''+b(x)u'=0,\quad x\in\Omega \\
u(0) &= u_0,\quad u(1)=u_1
\end{align}
\end{subequations}
with $b(x)$ bounded in $\Omega$
\end{frame}

\begin{frame} 
Let us approximate on $\bar{\omega}$ using \emph{centred differences} (a finite difference scheme):
\begin{subequations}
\begin{align}
y_0 &= u_0 \\
-D_+D_-y_i +b_iD_0y_i &= 0, \quad i=1,\ldots,n-1 \\
y_n &= u_1
\end{align}
\end{subequations}
where $b_i=b(x_i)$. The notation $D_+D_-$ and $D_0$ refer to \emph{difference operators}:
\begin{subequations}
\begin{align}
D_+y_i &= \frac{y_{i+1}-y_i}{h} \\
D_-y_i &= \frac{y_{i}-y_{i-1}}{h} \\
D_0y_i &= \frac{y_{i+1}-y_{i-1}}{2h} \\
D_+D_-y_i &= \frac{y_{i-1}-2y_i+y_{i+1}}{h^2}
\end{align}
\end{subequations}
These operators ``encode'' derivatives on the discrete grid used
\end{frame}


\begin{frame} 
For $i=1,\ldots,n-1$, define $\gamma_i = {hb_i}/2$. 
Then the problem can be written in matrix form as
\begin{equation}
Ay=f
\end{equation}
where $f=(u_0,0,\ldots,0,u_1)^T=u_0e_1+u_1e_{n+1}$ and
{\footnotesize
\[
A=
\begin{pmatrix}
1 & 0 & \\
-(1+\gamma_1) & 2 & -(1-\gamma+1) & \\
& -(1+\gamma_2) & 2 & -(1-\gamma_2) & \\
& & & \\
& & \ddots & \ddots & \ddots \\
& & & \\
& & & -(1+\gamma_{n-1}) & 2 & -(1-\gamma_{n-1}) \\
& & & & 0 & 1
\end{pmatrix}
\]
}

Properties of the (approximate) solution then depend on the properties of the M-matrix $A$
\end{frame}






%%%%%%%%%%%%%%%%%%
%%%%%%%%%%%%%%%%%%
%%%%%%%%%%%%%%%%%%
%%%%%%%%%%%%%%%%%%
\begin{frame}[allowframebreaks]
\frametitle{References}
\bibliographystyle{amsalpha}
\bibliography{MATH-4370-7370-lecture-notes.bib}
\end{frame}

\end{document}
